% Elsevier format paper
\documentclass[5p, times]{elsarticle}

% Required packages for Elsevier format
\usepackage[utf8]{inputenc}
\usepackage{lineno,hyperref}
\modulolinenumbers[5]

% Additional packages
\usepackage{amsmath,amssymb,amsfonts}
\usepackage{graphicx}
\usepackage{booktabs}
\usepackage{adjustbox}
\usepackage{subcaption}
\usepackage{xcolor}
\usepackage{listings}
\usepackage{array}
\usepackage{diagbox}
\usepackage{makecell}
\usepackage{algorithm}
\usepackage{algpseudocode}
\usepackage{tikz}
\usetikzlibrary{shapes.geometric, arrows.meta, positioning, fit, backgrounds}

% PACKAGES
% ========================
\usepackage[utf8]{inputenc}
\usepackage[T1]{fontenc}
\usepackage[english]{babel}
\usepackage{geometry}
\usepackage{lmodern}
\usepackage{graphicx}
\usepackage{amsmath}
\usepackage{amssymb}
\usepackage{booktabs}
\usepackage{multirow}
\usepackage{xcolor}
\usepackage{hyperref}
\usepackage{titlesec}
\usepackage{float}
\usepackage{caption}
\usepackage{subcaption}
\usepackage[protrusion=true,expansion=true,final,kerning=true,spacing=true]{microtype}

\usepackage{enumitem}
\usepackage{ragged2e}
\usepackage{adjustbox}
\usepackage{tabularx}
\usepackage{algorithm}
\usepackage{algpseudocode}
\usepackage{verbatim}
\usepackage{fancyvrb}


% Math operators
\DeclareMathOperator*{\argmin}{arg\,min}
\DeclareMathOperator*{\argmax}{arg\,max}

% Bibliography style
\bibliographystyle{unsrt}

% Listings configuration
\lstset{
  basicstyle=\ttfamily\footnotesize,
  breaklines=true,
  showstringspaces=false,
  frame=single,
  backgroundcolor=\color{gray!10},
  keywordstyle=\color{blue}\bfseries,
  commentstyle=\color{green!60!black},
  stringstyle=\color{red},
  numberstyle=\tiny\color{gray},
  captionpos=b,
}

\usepackage[inline]{trackchanges}
%finalold  - Reject all edits.
%finalnew  - Accept all edits.
%footnotes - Display edits as footnotes.
%margins   - Display edits as margin notes.
%inline    - Display edits inline.

% For strikethrough in remove/change
\usepackage[normalem]{ulem}
\usepackage{soul}

% Revision highlighting commands
\definecolor{addgreen}{rgb}{1, 1, 1}      % Light green background
\definecolor{removeRed}{rgb}{1, 0.85, 0.85}     % Light red background
\definecolor{changeYellow}{rgb}{1, 1, 0.85}    % Light yellow background

% Override trackchanges add command with soul highlight (supports line breaks)
\sethlcolor{addgreen}
\renewcommand{\add}[1]{\hl{#1}}
\renewcommand{\remove}[1]{\colorbox{removeRed}{\sout{#1}}}
\renewcommand{\change}[2]{\colorbox{removeRed}{\sout{#1}} $\rightarrow$ \colorbox{addgreen}{#2}}

% Note: For tables, use \colorbox{addgreen}{[NEW]} in caption instead of wrapping

% Document begins
\begin{document}

\begin{frontmatter}

% Title
\title{A Multimodal Root Cause Analysis Engine for Cloud and Microservice Systems}

% Authors
\author[inst1]{Fulchany Arvelin Nanitelamio}
\ead{nanifulchany@gmail.com}

\author[inst1]{Süleyman Eken\corref{cor1}}
\ead{suleyman.eken@kocaeli.edu.tr}

\cortext[cor1]{Corresponding author: Süleyman Eken (suleyman.eken@kocaeli.edu.tr)}

\author[inst2]{Ümit Demirbaga}
\ead{umitdemirbaga@ankara.edu.tr}

% Affiliations
\affiliation[inst1]{organization={Information Systems Engineering, Kocaeli University},
            addressline={Umuttepe Campus},
            city={Kocaeli},
            postcode={41001},
            country={Türkiye}}

\affiliation[inst2]{organization={Software Engineering, Ankara University},
            city={Ankara},
            postcode={06560},
            country={Türkiye}}

% Abstract
\begin{abstract}
Modern cloud computing infrastructures and microservice architectures have transformed system failure diagnosis (Root Cause Analysis, RCA) into an exceptionally complex problem. This paper presents four integrated contributions. First, a synchronized multimodal synthetic dataset (Log, Metric, Dashboard\,PNG) spanning 11 infrastructure layers, 22 anomaly scenarios, and 10,000 samples is introduced together with a Drain3-based feature pipeline; the resulting Random Forest achieves 98.45\,\% F1 on anomaly detection, 100\,\% on layer classification, and 99.10\,\% on scenario classification. Second, XGBoost and LightGBM gradient-boosting classifiers are added to the seven-model suite and benchmarked with SHAP (SHapley Additive exPlanations) explainability; XGBoost reaches F1\,=\,0.98 and LightGBM F1\,=\,0.99, with rigorous overfitting analysis confirming a minimal train-test gap of $\Delta$F1\,=\,0.0125. Third, a dual-engine LLM layer (Ollama/Llama\,3.2 locally; Google Gemini\,2.5 Flash via API) enriches ML detections with natural-language explanations, ranked remediation runbooks, and executive reports; five metrics (Response Latency, JSON Parse Success Rate, Semantic Coherence, Remediation Relevance, and Executive Report Quality) are benchmarked across both engines. Fourth, validation on three public RCA corpora (HDFS 575\,K sessions, BGL 4.7\,M lines, OpenStack 137\,K lines) confirms $\Delta$F1\,$<$\,+0.014 generalisation for all models, with LightGBM achieving the highest cross-dataset mean F1 of 0.99.
\end{abstract}

% Keywords
\begin{keyword}
Kubernetes \sep root cause analysis \sep multimodal Learning \sep cloud computing \sep anomaly detection \sep real-world validation
\end{keyword}

\end{frontmatter}

% Line numbering for review
\linenumbers

% Main content sections
\section{Introduction}

The migration from monolithic architectures to Kubernetes-based microservices over the past decade has accelerated software delivery while raising operational complexity to unprecedented levels. In a
modern cloud environment, a single user request may traverse dozens of services, databases, and network layers within milliseconds. In such cascading structures, a minor memory leak or a slow database query can propagate into a system-wide outage within minutes~\cite{soldani2022}. Site Reliability Engineers (SREs) responding to such incidents must simultaneously correlate alerts from metric dashboards, raw log streams, and distributed traces, creating an information-overload problem that cannot be addressed by manual inspection alone.

The classical observability triad rests on three data modalities: (1)~Metrics: numerical time series (CPU, memory, disk I/O, network throughput); (2)~Logs: text-based, timestamped records of discrete software and hardware events; and
(3)~Dashboards: visual monitoring interfaces (e.g., Grafana) that aggregate both modalities into operator-readable panels. Most existing RCA tools consume only one of these modalities, leaving
systematic blind spots: a CPU spike visible in metrics may be explained by an application exception visible only in logs, while the dashboard reveals the temporal correlation between the two~\cite{instantops2024}. This single-modality limitation becomes especially critical in infrastructure-level failure scenarios, where the physical or virtual layer responsible for a degradation (whether a misconfigured database, an exhausted Linux VM, or a saturated CDN edge node) can only be pinpointed by combining quantitative metric signals with the qualitative, event-specific information embedded in system logs.

A secondary challenge is explainability: production AIOps systems must not only detect anomalies but justify their detections in terms interpretable by human operators. Black-box models that achieve high F1 scores in offline evaluation are routinely rejected by operations teams if they cannot articulate why a specific alert was raised~\cite{zhang2019cross}. Regulatory frameworks in the financial and healthcare sectors increasingly mandate audit trails for automated decision-making in critical infrastructure, further reinforcing the demand for model-level feature attribution.

A third challenge is autonomous remediation: as
microservice deployments scale beyond the capacity of manual runbook execution, SREs require systems that can translate a detected anomaly directly into an ordered, verified sequence of shell commands and Kubernetes operations, reducing mean time to recovery (MTTR) from tens of minutes to single-digit minutes.

This paper presents the multimodal RCA Engine, a four-part research program that addresses all three challenges through an integrated pipeline combining gradient-boosted classifiers, SHAP explainability, a dual-engine LLM layer, and rigorous multi-dataset validation. We make the following contributions:

\begin{enumerate}[leftmargin=*, topsep=2pt, itemsep=2pt]
    \item A fully labelled, synchronized multimodal synthetic dataset (10,000 samples, 11 layers, 22 scenarios) together with a Drain3 and statistical-metric feature pipeline, achieving F1\,=\,100\,\% on layer classification. The dataset is designed to be strictly harder than real-world benchmarks through controlled noise injection and feature dropout, ensuring conservative generalisation bounds.

    \item Integration of XGBoost and LightGBM into the seven-model RCA suite, with SHAP TreeExplainer analysis revealing the multimodal feature hierarchy that drives detections and satisfying operator-facing explainability requirements without post-hoc approximation.

    \item A dual-engine LLM benchmarking framework with five semantic metrics comparing a local open-source model (Llama\,3.2) against a cloud API model (Gemini\,2.5 Flash) for autonomous anomaly explanation and self-remediation generation, including a structured fallback cascade that guarantees pipeline robustness under partial LLM failure.

    \item Comprehensive validation on three public log corpora (HDFS, BGL, OpenStack) replacing single-dataset evaluation and demonstrating cross-domain generalisation of the complete model suite, with LightGBM achieving the highest cross-dataset mean F1 of 0.99.
\end{enumerate}

The remainder of this paper is structured as follows. Section \ref{related} reviews related work across rule-based approaches, log analysis, gradient boosting, SHAP explainability, multimodal integration, and the LogHub corpora. Section \ref{sec3} describes the system
architecture and synthetic data generation pipeline. Section \ref{sec4} details feature extraction and multimodal fusion. Section \ref{sec5} presents the gradient boosting models and SHAP analysis. Section \ref{sec6} reports full model suite results on synthetic data. Section \ref{sec7} describes the dual-engine LLM benchmarking. Section \ref{sec:realdata} validates the pipeline on real-world datasets. Sections \ref{sec9} and \ref{sec10} discuss results and give conclusion and future directions, respectively.

% ========================
% 2. RELATED WORK
% ========================
\section{Related Work and Background}\label{related}

\subsection{Rule-based and Metric-based Approaches}
Early RCA tools relied on static threshold rules that become unmanageable in dynamic cloud environments. Metric-based anomaly detection progressively shifted from ARIMA and Holt-Winters seasonal decomposition toward unsupervised learning: Isolation Forest~\cite{soldani2022}, One-Class SVM, and LSTM-Autoencoders capture non-linear temporal patterns missed by statistical models. These approaches detect periodic deviations but remain semantically blind: they cannot distinguish a CPU surge caused by a memory leak from one caused by a legitimate traffic spike. Sharma et
al.~\cite{sharma2023microservice} observe that threshold-based systems generate excessive false positives in autoscaling environments because resource limits fluctuate dynamically, making
static thresholds structurally inappropriate. Recent works have therefore turned to learned, adaptive detection boundaries that incorporate contextual metadata such as deployment configuration
and service dependency topology~\cite{chen2020unsupervised}.

\subsection{Log Analysis with Drain and Derivatives}
Drain~\cite{drain2017} clusters raw log lines into fixed templates via a depth-bounded Parse Tree, masking variable tokens (IP addresses, block IDs, transaction identifiers) to produce stable,
reusable templates. Drain3\footnote{https://pypi.org/project/drain3/}, its actively maintained Python extension, generates session-level TF-IDF vectors and event-count matrices that serve as the log feature backbone throughout this work. On the HDFS, BGL, and OpenStack corpora from LogHub~\cite{loghub2023}, Drain3 achieves template parsing accuracy above 97\,\%, making it the
de-facto standard for production log pre-processing. Alternative parsers such as AEL~\cite{jiang2008abstracting},
IPLoM~\cite{makanju2009clustering}, and Spell~\cite{du2016spell} have been evaluated systematically in the LogPai benchmark~\cite{zhu2019tools}, confirming Drain's superior combination of accuracy and throughput at scale. More recent neural approaches such as LogBERT~\cite{guo2021logbert} achieve higher template accuracy on ambiguous log formats but incur significantly higher computational cost, making Drain3 the practical choice for real-time RCA pipelines with sub-second latency requirements.

\subsection{Gradient Boosting in AIOps}
XGBoost~\cite{xgboost2016} and LightGBM~\cite{lightgbm2017} have
become the default baseline in tabular anomaly detection due to their regularised tree ensembles, native handling of missing values, and fast inference. In the AIOps context, XGBoost has been applied to KPI anomaly scoring~\cite{zhang2019cross} and LightGBM to log-based failure classification~\cite{he2016evaluation}. Their
integration with multimodal (log and metric) features and subsequent SHAP interpretation in a hierarchical RCA pipeline, however, has not been systematically studied. Both methods outperform deep learning alternatives on tabular data with fewer than $10^5$ training samples~\cite{grinsztajn2022tree}, a regime directly applicable to the 10{,}000-sample synthetic dataset introduced in this work.

\subsection{SHAP Explainability}
SHAP~\cite{shap2017} attributes each prediction to individual input features using Shapley values from cooperative game theory. For tree-based models, the TreeSHAP variant achieves exact computation in polynomial time $O(TLD^2)$ by exploiting conditional independence in tree structures. In AIOps, SHAP explanations enable SRE engineers to understand why a detection was flagged, a critical requirement when operators must decide whether to escalate an alert or dismiss it as a false positive. Lundberg et al.~\cite{lundberg2020local2global} further demonstrate that global SHAP feature importance aggregated over a test set reveals the systematic decision-making logic of ensemble models, providing a
model-level audit trail suitable for compliance reporting. LIME~\cite{ribeiro2016lime}, an alternative local explanation framework, provides approximate rather than exact Shapley values and has lower fidelity for tree-based models; TreeSHAP is therefore preferred in this work.

\subsection{Multimodal Integration and LLM}
Recent works demonstrate that fusing logs, metrics, and traces outperforms single-modal baselines. AnoFusion~\cite{zhao2023anofusion} uses a Graph Transformer Network to learn heterogeneous multimodal
correlations, followed by GAT and GRU to detect anomalies at the system level. DiagFusion~\cite{zhang2023diagfusion} employs
fastText embeddings and a dependency graph GNN for root-cause typing. InstantOps~\cite{instantops2024}, the closest comparison point, combines GNN and GRU to model spatial and temporal dependencies in microservice call graphs, achieving F1\,=\,0.98 on
MicroSS and F1\,=\,0.96 on QoTD. On the LLM integration side, recent systems such as RCACopilot~\cite{chen2023rcacopilot} and AIOpsLab~\cite{chen2024aiopslab} demonstrate that language models can substantially accelerate the human-in-the-loop diagnosis cycle when provided with structured, ML-derived anomaly context. Our work
differs from these systems in four respects: (a)~it targets infrastructure-layer classification rather than node-level service localisation; (b)~it incorporates a visual modality
(Dashboard\,PNG); (c)~it uses a fully open and reproducible synthetic dataset; and (d)~it adds a dual-engine LLM layer with quantitative benchmarking for autonomous remediation generation.

\subsection{Real-world Log Corpora}
The LogHub benchmark~\cite{loghub2023} provides three datasets used for validation in this work. HDFS\,v1 contains 575{,}061 block sessions, 29 event types, and a 2.9\,\% anomaly rate,
originating from an Apache Hadoop Distributed File System deployment. Anomaly labels were hand-annotated by domain experts and verified by
cross-referencing block-level error propagation chains. BGL logs 4.7\,M lines from a BlueGene/L supercomputer with an 8.3\,\% anomaly rate dominated by hardware fault codes and MPI communication errors; the temporal density of faults is substantially higher than HDFS, making BGL a complementary stress test for models trained on
balanced synthetic data. OpenStack provides 137{,}K lines from Nova and Neutron cloud platform services with a 10.5\,\% anomaly rate driven by VM provisioning failures and authentication faults;
its structured log format (consistent timestamp, service tag, and request ID fields) makes it the most similar real-world dataset to the synthetic logs generated in this work.

Existing approaches suffer from at least one of the following limitations: (i)~single-modality input (metric-only or log-only); (ii)~lack of operator-facing explainability; (iii)~absence of autonomous remediation generation; (iv)~evaluation on a single
dataset. The Multimodal RCA Engine addresses all four limitations simultaneously through a principled combination of multimodal feature fusion, TreeSHAP attribution, dual-engine LLM enrichment, and three-corpus validation.

% ========================
% 3. SYSTEM ARCHITECTURE AND DATA GENERATION
% ========================
\section{System Architecture and Multimodal Data Generation}\label{sec3}

Figure~\ref{fig:pipeline_full} illustrates the overall architecture of the Multimodal RCA Engine. Three data sources (metrics, logs, and dashboards) feed into an AI processing core composed of gradient-boosting classifiers and an LLM enrichment layer. Outputs are autonomous RCA reports and self-healing action recommendations.

\begin{figure*}[htbp]
    \centering
    \adjustbox{max width=.75\linewidth}{%
        \includegraphics{04_pipeline_architecture.jpg}}
    \caption{Multimodal RCA Engine overall architecture}
    \label{fig:pipeline_full}
\end{figure*}

\subsection{Synthetic Dataset Generator}

Cloud providers do not release real failure logs due to privacy and competitive constraints. The dataset\_generator module synthesises 10,000 labelled samples across 11 infrastructure layers: CDN, Firewall, Proxy, Kubernetes Ingress, Kubernetes
Deployments, Application, Database, Storage, Network, Linux VM, and Linux Physical Host. Two anomaly scenarios are defined per layer, yielding 22 anomaly scenarios in total. The class split is
exactly 50\,\% normal / 50\,\% anomalous, with 908--911 samples per infrastructure layer ensuring complete class balance. Anomaly severity is distributed as 2,955 high, 1,363 critical, and
682 medium instances, reflecting realistic severity skew observed in production incident datasets. Each of the 22 anomaly scenarios contains 227--228 samples, maintaining strict class balance across all scenarios spanning all layers from CDN to Linux Physical Host. Table~\ref{tab:dataset_overview} summarises the key dataset parameters, and Table~\ref{tab:scenarios} enumerates the complete catalogue of the 22 anomaly scenarios together with their infrastructure layer, root-cause category, severity level, and injected anomaly type.

\begin{table}[htbp]
\centering
\caption{Synthetic dataset overview}
\label{tab:dataset_overview}
\adjustbox{max width=\columnwidth}{%
\begin{tabular}{@{}ll@{}}
\toprule
\textbf{Parameter}            & \textbf{Value} \\ \midrule
Total samples                 & 10,000         \\
Infrastructure layers         & 11             \\
Anomaly scenarios             & 22             \\
Normal / Anomaly split        & 50\,\% / 50\,\% \\
Samples per layer             & 908--911       \\
Samples per scenario          & 227--228       \\
Severity: High / Critical / Medium & 2,955 / 1,363 / 682 \\
Log formats simulated         & HDFS, BGL, OpenStack, Apache \\
Metric signals per sample     & 5--8 time-series channels \\ \bottomrule
\end{tabular}}
\end{table}

\begin{table*}[htbp]
\centering
\caption{The 22 anomaly scenarios across the 11 infrastructure layers (2 scenarios per layer).}
\label{tab:scenarios}
\adjustbox{max width=\linewidth}{%
\footnotesize
\begin{tabular}{@{}llllll@{}}
\toprule
\textbf{ID} & \textbf{Infrastructure Layer} & \textbf{Anomaly Scenario} & \textbf{Root-Cause Category} & \textbf{Severity} & \textbf{Anomaly Type} \\ \midrule
CDN\_001  & CDN                     & Sudden Traffic Spike            & network  & critical & spike           \\
CDN\_002  & CDN                     & Response Time Increase          & software & high     & gradual\_rise   \\
FW\_001   & Firewall                & Increased Blocked Connections   & network  & high     & spike           \\
FW\_002   & Firewall                & High Bandwidth Usage            & network  & high     & sustained\_high \\
PRX\_001  & Proxy                   & Error Rate Surge                & software & high     & gradual\_rise   \\
PRX\_002  & Proxy                   & Response Latency                & resource & medium   & gradual\_rise   \\
K8I\_001  & Kubernetes Ingress      & High Error Rate                 & config   & high     & spike           \\
K8I\_002  & Kubernetes Ingress      & Unbalanced Traffic Distribution & config   & medium   & oscillation     \\
K8D\_001  & Kubernetes Deployments  & Pod Restart Count Increase      & resource & high     & spike           \\
K8D\_002  & Kubernetes Deployments  & CPU Usage Increase              & resource & high     & sustained\_high \\
APP\_001  & Application             & Processing Time Increase        & software & high     & gradual\_rise   \\
APP\_002  & Application             & Error Rate Increase             & software & critical & spike           \\
DB\_001   & Database                & Slow Query Increase             & software & high     & gradual\_rise   \\
DB\_002   & Database                & Connection Count Increase       & network  & critical & spike           \\
STR\_001  & Storage                 & High Disk Usage                 & resource & high     & gradual\_rise   \\
STR\_002  & Storage                 & Latency Increase                & hardware & high     & spike           \\
NET\_001  & Network                 & Packet Loss Increase            & network  & high     & spike           \\
NET\_002  & Network                 & Excessive Latency               & network  & medium   & sustained\_high \\
VM\_001   & Linux VM                & CPU Usage Surge                 & resource & high     & sustained\_high \\
VM\_002   & Linux VM                & Memory Exhaustion               & resource & critical & gradual\_rise   \\
HOST\_001 & Linux Host              & CPU Temperature Increase        & hardware & critical & gradual\_rise   \\
HOST\_002 & Linux Host              & Disk Error Count Increase       & hardware & critical & step\_increase  \\ \bottomrule
\end{tabular}}
\end{table*}

\subsubsection{Time-series Metric Simulation}
The metric\_generator applies a three-layer statistical model. The base signal combines a layer-specific mean and variance with a diurnal oscillation of period $T = 1{,}440$\,s and first-order autocorrelation $\rho = 0.7$. The noise layer
superimposes i.i.d.\ Gaussian noise $\mathcal{N}(0,\sigma^2)$ and random micro-spikes drawn from a Poisson process. The anomaly injection selects one of three modes:
\begin{align}
    \text{sudden\_spike:}\quad & x_t = \mu + k\sigma,\quad
        k \sim \mathcal{U}(3,8) \label{eq:spike}\\
    \text{gradual\_rise:}\quad & x_t = \mu +
        \frac{t - t_0}{T_a}\,\Delta \label{eq:rise}\\
    \text{oscillation:}\quad   & x_t = \mu +
        A\sin\!\left(\tfrac{2\pi t}{P}\right) \label{eq:osc}
\end{align}
where $t_0$ is the anomaly onset, $T_a$ the escalation window, $\Delta$ the target deviation amplitude, $A$ the oscillation amplitude, and $P$ its period. Table~\ref{tab:metric_examples}
presents representative generated metric values for normal and anomalous CDN sessions, illustrating the statistical divergence introduced by each anomaly mode.

\begin{table}[htbp]
\centering
\caption{Example of generated metric features: CDN layer (normal vs.\ anomaly)}
\label{tab:metric_examples}
\adjustbox{max width=\columnwidth}{%
\begin{tabular}{@{}lcccc@{}}
\toprule
\textbf{Feature} & \textbf{Normal} & \textbf{Spike} &
    \textbf{Gradual} & \textbf{Osc.} \\ \midrule
Mean ($\bar{x}$)      & 42.3  & 91.7  & 68.4  & 45.1  \\
Std dev ($\sigma_x$)  & 3.1   & 18.4  & 9.7   & 12.8  \\
Slope ($\beta$)       & 0.002 & 0.041 & 0.187 & 0.003 \\
Spike score           & 0.02  & 0.41  & 0.18  & 0.22  \\
Autocorrelation       & 0.71  & 0.43  & 0.68  & 0.55  \\
Volatility ($\nu$)    & 2.4   & 14.2  & 7.8   & 10.1  \\ \bottomrule
\end{tabular}}
\end{table}

\subsubsection{Log Session Generation}
The log\_generator produces logs faithful to HDFS, BGL, OpenStack, and Apache formats. During anomaly periods, log levels dynamically shift toward ERROR and CRITICAL: the fraction of CRITICAL lines increases from $\sim$4\,\% in normal sessions to
$\sim$21\,\% in anomalous ones. MAC addresses, IP addresses, and session IDs are randomly generated per sample to prevent spurious memorisation. Table~\ref{tab:log_examples} provides representative
log lines generated for normal and anomalous sessions across two infrastructure layers.

\begin{table}[htbp]
\centering
\caption{Example of generated log lines: normal vs.\ anomaly}
\label{tab:log_examples}
\small
\begin{tabularx}{\columnwidth}{@{}l l X@{}}
\toprule
\textbf{Layer} & \textbf{Level} & \textbf{Log Message Template} \\ \midrule
CDN      & INFO   & Cache hit for block \texttt{<ID>} from \texttt{<IP>} \\
CDN      & WARN   & Cache miss rate exceeding threshold on \texttt{<ID>} \\
CDN      & ERROR  & Response timeout \texttt{<N>}\,ms for request \texttt{<ID>} \\
CDN      & CRIT.  & Edge node \texttt{<ID>} unreachable; failover triggered \\
\midrule
Linux VM & INFO   & CPU usage \texttt{<N>}\,\% on vm-\texttt{<ID>} \\
Linux VM & WARN   & Memory pressure: \texttt{<N>}\,MB available \\
Linux VM & ERROR  & OOM-killer invoked; process \texttt{<N>} terminated \\
Linux VM & CRIT.  & Kernel panic on vm-\texttt{<ID>}: out of memory \\ \bottomrule
\end{tabularx}
\end{table}

\subsubsection{Dashboard Image Generation}
The dashboard\_generator renders Grafana dark-theme PNG images. A red vertical line marks the anomaly onset; the anomaly interval is highlighted with a translucent red background. As shown in Figure~\ref{fig:dashboards}, the four dashboard examples cover CDN normal operation, CDN Sudden Traffic Spike, Database Slow Query Increase, and Linux VM CPU Usage Surge, illustrating the visual
discriminability of anomaly patterns across layers. This visual modality is currently used for qualitative validation; full VLM integration is planned as future work.

\begin{figure*}[htbp]
    \centering
    \adjustbox{max width=\linewidth}{%
        \includegraphics{01_dataset_exploration_figure_6.png}}
    \caption{Generated dashboard image examples (Grafana dark theme).
    \textit{Top-left:} CDN normal operation. \textit{Top-right:}
    CDN: Sudden Traffic Spike (CRITICAL). \textit{Bottom-left:}
    Database: Slow Query Increase (HIGH). \textit{Bottom-right:}
    Linux VM: CPU Usage Surge (HIGH). Red vertical line: anomaly
    onset; red background: anomaly window.}
    \label{fig:dashboards}
\end{figure*}

\subsection{Root Cause Taxonomy}
Each of the 22 anomaly scenarios is assigned to exactly one of five
root-cause categories. Because a single infrastructure layer hosts
two distinct scenarios, a layer may span more than one category;
for example, the Database layer contributes a \textit{software}
Slow Query Increase scenario and a \textit{network} Connection
Count Increase scenario. Aggregated over all 22 scenarios, the
category distribution is: \textit{network} (27.3\,\%, 6 scenarios:
CDN, Firewall $\times$2, Database, Network $\times$2);
\textit{resource} (27.3\,\%, 6 scenarios: Proxy, Kubernetes
Deployments $\times$2, Storage, Linux VM $\times$2);
\textit{software} (22.7\,\%, 5 scenarios: CDN, Proxy, Application
$\times$2, Database); \textit{hardware} (13.6\,\%, 3 scenarios:
Storage, Linux Host $\times$2); and \textit{config} (9.1\,\%,
2 scenarios: Kubernetes Ingress $\times$2). The complete
per-scenario assignment is listed in Table~\ref{tab:scenarios},
and the cross-matrix of layer vs.\ root-cause category is
illustrated in Figure~\ref{fig:root_cause_matrix}.

\begin{figure*}[htbp]
    \centering
    \adjustbox{max width=\linewidth}{%
        \includegraphics{01_dataset_exploration_figure_7.png}}
    \caption{\textit{Left:} Root-cause category distribution
    \textit{Right:} Layer $\times$ Root-Cause cross-matrix}
    \label{fig:root_cause_matrix}
\end{figure*}

% ========================
% 4. FEATURE EXTRACTION
% ========================
\section{Feature Extraction and Multimodal Fusion}\label{sec4}

\subsection{Log Feature Extraction}
Drain3's Parse Tree groups log lines by shared token structure, masking variable fields (IPs, block IDs, timestamps) to produce stable templates. Two feature families are derived per session:
\begin{enumerate}[leftmargin=*, topsep=2pt, itemsep=1pt]
    \item TF-IDF vectors (100 dimensions): importance-weighted
    bag-of-templates encoding which log event types appear and how
    frequently relative to the corpus.
    \item Event-count matrices: per-session raw frequencies
    of each template and of each log level (INFO, WARN, ERROR,
    CRITICAL).
\end{enumerate}

Applied to the HDFS dataset, the top-15 templates extracted by Drain3 include ``Receiving block'', ``Received block'', and ``PacketResponder'', which carry the highest anomaly signal frequency. The discriminative power of log features is illustrated
by the following statistics computed on the synthetic dataset: the mean ERROR line count rises from 5.7 (normal) to 14.2 (anomalous),
a +149\,\% increase; the CRITICAL line count rises from 1.7 to 8.5, a +400\,\% increase. As shown in
Figure~\ref{fig:log_levels}, the dramatic shift in ERROR and CRITICAL counts between normal and anomalous sessions confirms the critical contribution of log features to anomaly detection.

\begin{figure*}[h!]
    \centering
    \adjustbox{max width=0.75\linewidth}{%
        \includegraphics{01_dataset_exploration_figure_3.png}}
    \caption{Mean log-level distribution in normal vs.\ anomalous sessions}
    \label{fig:log_levels}
\end{figure*}

\subsection{Metric Feature Extraction}
Raw time-series windows are transformed into the following scalar feature vector via a sliding-window computation:
\begin{equation}
    f_{\text{metric}} = \bigl[\bar{x},\;\tilde{x},\;
    \sigma_x,\;\beta,\;\nu,\;s\bigr]
\end{equation}
where $\bar{x}$ is the sample mean, $\tilde{x}$ the median, $\sigma_x$ the standard deviation, $\beta$ the linear slope (estimated by ordinary least squares over the window), $\nu$ the volatility (mean absolute deviation from $\bar{x}$), and $s$ the
spike score (fraction of samples exceeding $\bar{x} + 2\sigma_x$). To prevent data leakage, metric column names are masked as m0, m1, \ldots\ Figure~\ref{fig:metric_timeseries} illustrates the CDN layer metric time series, contrasting normal
diurnal oscillation with the sharp degradation signature of a Sudden Traffic Spike anomaly.

\begin{figure*}[h!]
    \centering
    \adjustbox{max width=\linewidth}{%
        \includegraphics{01_dataset_exploration_figure_4.png}}
    \caption{CDN layer metric time series. \textit{Top:} Normal
    operation: metrics exhibit regular diurnal oscillation.
    \textit{Bottom:} Anomaly (Sudden Traffic Spike): from
    $t=13$\,min (red vertical line) Error Rate and Response Time
    spike sharply while Cache Hit Ratio collapses. Red background
    marks the anomaly window.}
    \label{fig:metric_timeseries}
\end{figure*}

\subsection{Unified Feature Matrix and Data Augmentation}
Dense metric features and sparse log TF-IDF features are horizontally concatenated:
\begin{equation}
    X = \bigl[X_{metric} \;\|\;
    X_{tfidf} \;\|\;
    X_{event-count}\bigr]
\end{equation}
The MultimodalDataLoader further injects Gaussian noise ($\sigma_{noise} = 0.25$) and random feature dropout ($p_{drop} = 0.08$) to simulate real-world sensor degradation and missing telemetry, intentionally making the synthetic training distribution harder than the real datasets.

\subsection{Hierarchical Classification Strategy}
The detection pipeline operates in two sequential phases:
\begin{enumerate}[leftmargin=*, topsep=2pt, itemsep=1pt]
    \item Phase 1 -- Binary Anomaly Detector: classifies each observation as Normal\,(0) or Anomalous\,(1).
    \item Phase 2 -- RCA Multi-Class Classifiers: applied only to observations flagged as anomalous. Four parallel classifiers predict: Infrastructure Layer (11 classes), Anomaly
    Scenario (22 classes), Root-Cause Category (5 classes), and Fault Severity (3 classes).
\end{enumerate}

% ========================
% 5. GRADIENT BOOSTING MODELS AND SHAP
% ========================
\section{Gradient Boosting Models and SHAP Explainability}\label{sec5}

\subsection{XGBoost: Regularised Gradient Boosting}
XGBoost~\cite{xgboost2016} minimises a regularised objective at each boosting round $t$. Using a second-order Taylor expansion of the loss $\ell$, the objective for adding tree $f_t$ becomes:
\begin{equation}
    \mathcal{L}^{(t)} = \sum_{i=1}^{n}
    \bigl[g_i f_t(\mathbf{x}_i) +
    \tfrac{1}{2}h_i f_t^2(\mathbf{x}_i)\bigr] + \Omega(f_t)
\end{equation}
where $g_i = \partial_{\hat{y}^{(t-1)}}\ell(y_i,\hat{y}^{(t-1)})$
and $h_i = \partial^2_{\hat{y}^{(t-1)}}\ell(y_i,\hat{y}^{(t-1)})$
are the first and second gradient statistics, and the regularisation term penalises tree complexity:
\begin{equation}
    \Omega(f_t) = \gamma T + \tfrac{1}{2}\lambda\|\mathbf{w}\|^2
\end{equation}
with $T$ the number of leaves, $\mathbf{w}$ the leaf weight vector, $\gamma$ the minimum gain required for a split (pruning threshold), and $\lambda$ the L2 ridge penalty. The closed-form optimal weight
for leaf $j$ is:
\begin{equation}
    w_j^* = -\frac{\sum_{i\in I_j}g_i}{\sum_{i\in I_j}h_i + \lambda}
\end{equation}
Configuration: $n_{\text{est}}=200$, $\text{max\_depth}=8$,
$\eta=0.1$, subsample\,$=0.8$, colsample\,$=0.8$.

\subsection{LightGBM: GOSS and EFB}
LightGBM~\cite{lightgbm2017} shares the regularised boosting objective but introduces two key algorithmic departures.

Gradient-based One-Side Sampling (GOSS) retains all
instances with large gradients (top fraction $a$) and subsamples a fraction $b$ of those with small gradients, reweighting them by $\frac{1-a}{b}$ to preserve the gradient distribution.
Let $c = \frac{1-a}{b}$,
$G_l = \sum_{x_i\in A_l}g_i - c\sum_{x_i\in B_l}g_i$, and
$G_r = \sum_{x_i\in A_r}g_i - c\sum_{x_i\in B_r}g_i$.
The variance gain estimate for split $d$ on feature $j$ is:
\begin{equation}
    \tilde{V}_j(d) = \frac{1}{n}\!\left[
    \frac{G_l^2}{n_l^A + c\,n_l^B}
    +
    \frac{G_r^2}{n_r^A + c\,n_r^B}
    \right]
\end{equation}
where $A_l,A_r$ (resp.\ $B_l,B_r$) denote the large-gradient
(resp.\ small-gradient) instances falling left and right of
split $d$.

Exclusive Feature Bundling (EFB) merges mutually exclusive sparse features (those that rarely take non-zero values simultaneously) into a single dense bundle, reducing the effective feature dimensionality from $d$ to $d'<d$ without information loss. Together, GOSS and EFB reduce histogram-building complexity from $O(n\cdot d)$ to $O(n'\cdot d')$ where $n' \ll n$, enabling
significantly faster training on high-dimensional feature sets.

\subsection{XGBoost vs.\ LightGBM Performance}

Table~\ref{tab:gb_perf} reports the head-to-head performance of both gradient-boosting models on the synthetic test set. LightGBM outperforms XGBoost by $\Delta$F1\,=\,+0.003, attributable to GOSS's more efficient exploitation of informative gradient instances in the balanced 50/50 class distribution. Both models achieve AUC\,=\,0.99, confirming near-perfect probabilistic ranking.

\begin{table}[h!]
\centering
\caption{Gradient boosting performance on synthetic test set}
\label{tab:gb_perf}
\adjustbox{max width=\columnwidth}{%
\begin{tabular}{@{}lccccc@{}}
\toprule
\textbf{Model} & \textbf{Acc.} & \textbf{Prec.} & \textbf{Rec.}
    & \textbf{F1} & \textbf{AUC} \\ \midrule
XGBoost  & 0.98 & 0.9880 & 0.9870 & 0.9875 & 0.9994 \\
LightGBM & 0.9905 & 0.9910 & 0.9900
    & 0.9905 & 0.9994 \\ \bottomrule
\end{tabular}}
\end{table}

\subsection{Overfitting Analysis}

\subsubsection{Train / Validation / Test Comparison}
A 70/10/20 stratified split yields for XGBoost: Train F1\,=\,1.0000, Validation F1\,=\,0.9880, Test F1\,=\,0.9875. The train-test gap $\Delta$F1\,=\,0.0125 is intentionally induced by the data loader's noise injection (noise\_level\,=\,0.15,
feature\_dropout\,=\,0.05).

\subsubsection{5-Fold Stratified Cross-Validation}
Running Stratified K-Fold CV on the full dataset:
\begin{equation}
    \overline{F1}_{CV} = 0.9895 \pm 0.0016
\end{equation}
The low standard deviation ($\sigma<0.005$) confirms stable performance across folds, a hallmark of good generalisation rather than memorisation.

\subsubsection{Regularised XGBoost}
A second XGBoost variant is trained with stronger regularisation: $max\_depth=5$, $\gamma=1.0$, L1\,$\alpha=0.1$,
L2\,$\lambda=2.0$, $min\_child\_weight=5$, with early
stopping at round 194. The regularised model reduces the train-test gap while maintaining test F1 within $\pm0.003$ of the base model, confirming that the original configuration is not over-fit in a
clinically significant sense.

\subsection{SHAP Explainability Analysis}

SHAP TreeExplainer computes exact Shapley values for the XGBoost
detector on 500 held-out samples. For a model with $M$ features,
the Shapley value of feature $j$ for prediction $i$ is:
\begin{equation}
    \phi_j^{(i)} = \sum_{S \subseteq F \setminus \{j\}}
    \frac{|S|!\,(M-|S|-1)!}{M!}
    \bigl[v_i(S\cup\{j\}) - v_i(S)\bigr]
\end{equation}
where $v_i(S)$ is the model output conditioned on the feature
coalition $S$. TreeSHAP evaluates this in $O(TLD^2)$ time using
the tree's conditional independence structure.

\subsubsection{Global Feature Importance}
As shown in Figure~\ref{fig:shap_importance}, the SHAP summary reveals a three-tier hierarchy:

\begin{enumerate}[leftmargin=*, topsep=2pt, itemsep=1pt]
    \item Log features (log\_severity\_score, log\_error\_ratio, log\_critical\_ratio) dominate the top positions. Sharp level transitions at anomaly onset generate the strongest, most immediate signal.
    \item TF-IDF features (tfidf\_37,
    tfidf\_26) appear at mid-rank, capturing template
    frequency anomalies invisible in the metric stream.
    \item Metric features (m1\_autocorr, m0\_slope, m2\_spike\_score) contribute in lower positions but provide the earliest signal for gradual-rise scenarios~\eqref{eq:rise}, where log levels change only after metrics have already elevated.
\end{enumerate}

This hierarchy directly validates the multimodal design: a metrics-only system cannot achieve 100\,\% layer accuracy because the same CPU spike can originate from multiple layers; only the log template disambiguates the source.

\begin{figure*}[h!]
    \centering
    \adjustbox{max width=0.90\linewidth}{%
        \includegraphics{02_model_training_figure_2.png}}
    \caption{Top-20 SHAP feature importance (XGBoost, TreeSHAP).
    Log features dominate, followed by TF-IDF and metric
    autocorrelation features, validating the multimodal fusion
    design.}
    \label{fig:shap_importance}
\end{figure*}

\subsubsection{Individual Anomaly Explanation}
For scenario FW\_002 (High Bandwidth Utilisation), the SHAP waterfall plot decomposes the prediction from the base value $\mathbb{E}[f(\mathbf{x})]$ to the final output. Large positive
contributions from log\_error\_ratio and tfidf\_37
drive the prediction toward ``anomaly'', while m0\_mean provides a small negative correction reflecting that some metric values remain in the normal range, illustrating how contradictory
signals across modalities are resolved by the gradient-boosted ensemble.

% ========================
% 6. MODEL SUITE AND SYNTHETIC EXPERIMENTS
% ========================
\section{Full Model Suite and Synthetic Data Experiments}\label{sec6}

The 10,000-sample dataset is split 80/20 (train/test) with Stratified 5-Fold Cross-Validation.

\subsection{Phase 1: Binary Anomaly Detection}

Table~\ref{tab:binary_perf} reports Binary Anomaly Detection performance. The Random Forest achieves F1\,=\,98.45\,\% with AUC-ROC\,=\,0.9988 and a confusion matrix of 986 true negatives, 983 true positives, 14 false positives, and 17 false negatives. This highly favourable operating point reflects the strong discriminability of the fused log\,+\,metric feature space.

\begin{table}[h!]
\centering
\caption{Performance results for binary anomaly detection }
\label{tab:binary_perf}
\adjustbox{max width=\columnwidth}{%
\begin{tabular}{@{}ll@{}}
\toprule
\textbf{Metric}                & \textbf{Value}     \\ \midrule
Accuracy                       & 98.45\,\%          \\
Precision                      & 98.49\,\%          \\
Recall                         & 98.40\,\%          \\
F1 Score                       & 98.45\,\%\\
AUC-ROC                        & 0.9988             \\
Confusion Matrix (TN/FP/FN/TP) & 986/14/17/983      \\ \bottomrule
\end{tabular}}
\end{table}

Figure~\ref{fig:binary_metrics} shows the confusion matrix, ROC curve (AUC=0.9988), and Precision-Recall curve (F1\,=\,0.9845) for the Random Forest anomaly detector, confirming the near-perfect
separation between normal and anomalous samples.

\begin{figure*}[h!]
    \centering
    \adjustbox{max width=\linewidth}{%
        \includegraphics{02_model_training_figure_1.png}}
    \caption{Random Forest anomaly detector performance.
    \textit{Left:} Confusion Matrix (986 TN, 983 TP, 14 FP, 17 FN).
    \textit{Centre:} ROC curve, AUC\,=\,0.9988.
    \textit{Right:} Precision-Recall curve, F1\,=\,0.9845.}
    \label{fig:binary_metrics}
\end{figure*}

\subsection{Phase 2: RCA Multi-Class Classification}

Table~\ref{tab:rca_perf} summarizes the full RCA pipeline performance across all four classification targets. Infrastructure Layer classification achieves a perfect F1\,=\,100.00\,\%, confirming complete metric/log separability across the 11
infrastructure layers. Root-Cause Category classification reaches F1\,=\,99.90\,\% with only one misclassification (a Linux VM CPU Surge partially overlapping with a network congestion metric pattern). Fault Severity and Anomaly Scenario classification achieve F1\,=\,99.80\,\% and 99.10\,\% respectively, with the scenario task being the most challenging due to within-layer scenario
similarity.

\begin{table}[h!]
\centering
\caption{Performance results for full RCA pipeline }
\label{tab:rca_perf}
\adjustbox{max width=\columnwidth}{%
\begin{tabular}{@{}lcc@{}}
\toprule
\textbf{Classification (\# of classes)}   & \textbf{Accuracy}    & \textbf{F1(w)} \\ \midrule
Binary (anomaly or not) (2)        & 98.45\,\%            & 98.45\,\%      \\
Root Cause Category (5)     & 99.90\,\%            & 99.90\,\%      \\
Fault Severity (3)          & 99.80\,\%            & 99.80\,\%      \\
Infrastructure Layer (11)   & 100.00\,\%  & 100.00\,\% \\
Anomaly Scenario (22)       & 99.10\,\%            & 99.10\,\%      \\ \bottomrule
\end{tabular}}
\end{table}

The single root-cause misclassification (resource $\to$ network) occurs for a Linux VM CPU Surge scenario whose metric signature partially overlaps with a network congestion pattern. This is the only case in which the multimodal feature vector is insufficient to break the tie, and it underscores the value of the dashboard visual modality as a third disambiguation channel for future work.

\subsection{Full 7-Model Cross-Validation Comparison}

Table~\ref{tab:cv_synth} presents the 5-Fold CV F1 scores for all seven models. MLP leads at F1\,=\,0.9915\,$\pm$\,0.0021, followed by LightGBM (0.9905\,$\pm$\,0.0014) and XGBoost (0.9895\,$\pm$\,0.0016). All models cluster in the narrow 0.986--0.9915 band, demonstrating the high quality of the synthetic dataset and confirming that even the simplest baseline (Logistic Regression, F1\,=\,0.9864) achieves near-ceiling performance.
The low standard deviations across all models ($\sigma < 0.004$) confirm stable, fold-independent generalisation. All models achieve AUC\,$>$\,0.998 on ROC analysis, indicating near-perfect probabilistic ranking regardless of classification threshold.

\begin{table}[h!]
\centering
\caption{5-Fold CV F1 comparison for all models}
\label{tab:cv_synth}
\adjustbox{max width=\columnwidth}{%
\begin{tabular}{@{}lcc@{}}
\toprule
\textbf{Model}              & \textbf{F1 (mean)} & \textbf{($\pm$ std)} \\ \midrule
MLP Neural Network          & 0.9915    & $\pm$0.0021 \\
LightGBM                    & 0.9905             & $\pm$0.0014 \\
XGBoost                     & 0.9895             & $\pm$0.0016 \\
Gradient Boosting (sklearn) & 0.9882             & $\pm$0.0018 \\
KNN ($k=5$)                 & 0.9876             & $\pm$0.0024 \\
Random Forest               & 0.9870             & $\pm$0.0019 \\
Logistic Regression         & 0.9864             & $\pm$0.0031 \\ \bottomrule
\end{tabular}}
\end{table}

% ========================
% 7. LLM INTEGRATION AND BENCHMARKING
% ========================
\section{LLM Integration and Dual-Engine Benchmarking}\label{sec7}

\subsection{Dual-Engine Architecture}
The llm\_engine module supports two backends deployed in a fallback cascade:
\begin{itemize}[leftmargin=*, topsep=2pt, itemsep=1pt]
    \item Local Engine (Ollama/Llama\,3.2): open-source model running on-premises. Guarantees data privacy for enterprise deployments with air-gapped networks or strict data residency requirements. Inference latency depends on local
    GPU/CPU resources.
    \item Cloud Engine (Google Gemini\,2.5 Flash): stronger reasoning capacity with cloud-scale training. Applied to anomaly explanation, extended remediation planning, log interpretation, Q\&A chatbot, and executive report generation.
\end{itemize}

\subsection{Prompt Engineering and Structured Output}
Each ML detection generates a structured RCA context object containing the detected layer, scenario ID, severity level, root-cause category, and log feature vector. This context is serialised into a structured prompt instructing the LLM to return a JSON object with four mandatory fields: explanation (natural-language description for SRE operators), remediation\_steps (ordered list of commands
and verification criteria), estimated\_resolution\_time (SLA-aware estimate), and confidence\_score (LLM self-assessed certainty). This design decouples the ML detection stage from the language generation stage while ensuring machine-parseable outputs compatible with automated runbook systems.

\subsection{Resilient Fallback Parsing}
Smaller local models frequently fail to produce valid JSON. A three-level fallback cascade handles this:
\begin{enumerate}[leftmargin=*, topsep=2pt, itemsep=1pt]
    \item Direct parse: json.loads() on the raw
    response string.
    \item Regex extraction: locate the first balanced
    $\{\ldots\}$ block via a bracket-counting automaton and
    re-attempt parsing.
    \item Keyword heuristics: scan for ANOMALY/
    NORMAL labels and numeric confidence patterns to
    reconstruct a minimal valid JSON object.
\end{enumerate}

\subsection{Adaptive Rate-Limit Backoff}
Gemini Flash API rate-limit responses (HTTP 429) trigger a 13-second exponential backoff on standard calls, extending to 65 seconds after repeated consecutive failures. This strategy avoids quota exhaustion while keeping median pipeline latency
below 3 seconds per detection.

\subsection{LLM Performance Metrics}
We define five metrics to evaluate LLM performance on the RCA enrichment task, inspired by the MFR/MAR RCA evaluation framework of InstantOps~\cite{instantops2024} and adapted to the language generation setting:
\begin{enumerate}[leftmargin=*, topsep=2pt, itemsep=1pt]
    \item Response latency ($\tau$, seconds): wall-clock time from prompt submission to complete token stream.
    \item JSON parse success rate ($\rho_{json}$): fraction of responses successfully parsed on the first fallback
    level (direct json.loads).
    \item Semantic coherence ($\kappa$, 1--5): expert
    rating on whether the explanation correctly references the flagged layer and scenario.
    \item Remediation relevance ($\rho_{rem}$, 1--5): expert rating on correctness and specificity of proposed remediation steps.
    \item Executive report quality ($\rho_{exec}$,
    1--5): expert rating on whether the executive summary correctly prioritises anomalies by severity, assigns responsibility, and provides a clear action plan with estimated resolution timelines.
\end{enumerate}

As reported in Table~\ref{tab:llm_bench}, Gemini\,2.5 Flash outperforms Llama\,3.2 across all five metrics. The most critical gap is in $\rho_{json}$ (0.97 vs.\ 0.71), which directly
affects pipeline robustness: a 29\,\% parse failure rate in production would require the full three-level fallback on nearly one in three calls. Semantic coherence (4.6 vs.\ 3.2) and remediation relevance (4.5 vs.\ 3.0) also show substantial
improvements, reflecting Gemini's superior instruction following on structured RCA prompts. Response latency is paradoxically lower for the cloud model (1.8\,s vs.\ 4.2\,s for Llama\,3.2 running on consumer-grade CPU), due to Gemini's dedicated inference infrastructure. Llama\,3.2 compensates with zero API cost, zero network egress, and full data locality, making it the preferred
choice for privacy-sensitive or regulated deployments.

\begin{table}[h!]
\centering
\caption{Results for dual-engine LLM benchmarking}
\label{tab:llm_bench}
\begin{tabular}{@{}lrrrrr@{}}
\toprule
\textbf{Engine} & $\tau$ & $\rho_{json}$ & $\kappa$ & $\rho_{rem}$ & $\rho_{exec}$ \\ \midrule
Llama\,3.2 (local)      & 4.2 & 0.71 & 3.2 & 3.0 & 2.8 \\
Gemini\,2.5 Flash (API) & 1.8 & 0.97 & 4.6
    & 4.5 & 4.7 \\ \bottomrule
\end{tabular}
\end{table}

\subsection{Self-Remediation Runbook Generation: Case Study VM\_002}

To illustrate the LLM layer's practical output, we present the complete condensed runbook generated by Gemini\,2.5 Flash for scenario VM\_002 (Critical: Memory Exhaustion on a Linux VM), alongside the verification criteria it autonomously derives.

\subsubsection{Phase 1: Memory Cleanup and Leak Isolation}
Gemini\,2.5 Flash structures the intervention around three ordered steps. Step A identifies the leaking process by ranking all running processes by resident memory:
\begin{verbatim}
ps aux --sort=-%mem | head -n 11
kubectl top pods -A --sort-by=memory
\end{verbatim}
Step B performs a safe kernel-buffer flush. The model
correctly prescribes a sync before cache eviction to
guarantee write-back of dirty pages, preventing data inconsistency:
\begin{verbatim}
sync
echo 3 > /proc/sys/vm/drop_caches
\end{verbatim}
Step C restarts the leaking service with a 30-second
grace period (allowing in-flight requests to drain) or evicts the offending Kubernetes pod:
\begin{verbatim}
systemctl restart <service>.service
kubectl delete pod <pod> -n <ns> \
  --grace-period=30
\end{verbatim}
Verification: \texttt{free -h} ``available'' must exceed 20\,\% of total RAM; kubectl describe node
MemoryPressure must return False.

\subsubsection{Phase 2: Swap Optimisation and Disk-Thrashing Prevention}
Once physical memory is partially reclaimed, Gemini\,2.5 Flash addresses swap over-use. Live paging activity is monitored via:
\begin{verbatim}
vmstat 1 10
\end{verbatim}
It then applies an SRE-standard \texttt{vm.swappiness} reduction
from the kernel default (60) to 10, forcing the kernel to exhaust
physical RAM before paging to disk:
\begin{verbatim}
sysctl vm.swappiness=10
echo "vm.swappiness=10" >> /etc/sysctl.conf
\end{verbatim}
If \texttt{free -m} confirms sufficient reclaimed RAM, the model prescribes a safe swap recycle:
\begin{verbatim}
swapoff -a && swapon -a
\end{verbatim}
Verification: swap usage below 10\,\% in
\texttt{swapon --show}; \texttt{vmstat} \texttt{si}/\texttt{so}
columns at or near zero.

\subsubsection{Risk and Scope Assessment}
Gemini\,2.5 Flash autonomously annotates each remediation phase with its risk reduction rationale: Phase~1 eliminates the Kernel
OOM-Killer risk of randomly terminating critical system daemons (\texttt{sshd}, \texttt{kubelet}); Phase~2 prevents disk-thrashing latency degradation from microsecond to millisecond/second range. Estimated resolution time before OOM escalation:
5--10 minutes. The same engine aggregates the full
10-anomaly batch into an executive report identifying 4 Critical and 6 High priority events, concludes that 83.3\,\% of monitored components are degraded, and outputs a three-action priority
sequence for immediate operator execution.

% ========================
% 8. REAL-WORLD DATASET VALIDATION
% ========================
\section{Real-world Dataset Validation}
\label{sec:realdata}

\subsection{HDFS LogHub Dataset}

\subsubsection{Dataset and Preprocessing}
HDFS\,v1 contains 575{,}061 block sessions parsed by Drain3 into a 29-event frequency matrix. The dataset exhibits a 97.1\,\% Success / 2.9\,\% Fail label distribution (16{,}838 anomalous sessions). Among the 29 event types, E5, E2, and E11 are the most
frequent; certain event types diverge dramatically between normal and anomalous profiles, providing strong template-level discriminative signal. The 2.9\,\% anomaly rate creates significant class imbalance managed via: (i)~stratified sampling to 30,000 sessions (3{,}000 anomalous + 27{,}000 normal); (ii)~\texttt{class\_weight='balanced'} for tree-based models; (iii)~StandardScaler normalisation; (iv)~Stratified 5-Fold CV.

\subsubsection{Performance Results}

As reported in Table~\ref{tab:real_test}, all seven models achieve F1\,$\geq$\,0.9950 on the HDFS test set, with LightGBM leading at F1\,=\,0.9984 and AUC\,=\,1.0000. Random Forest and Gradient Boosting both achieve only 2 total errors (1 FP\,+\,1 FN each), demonstrating that the block-level event profile of HDFS is highly discriminative for all ensemble methods. Logistic Regression shows
the highest error count (5 FP\,+\,1 FN), reflecting the non-linearity of HDFS event co-occurrence patterns that a linear boundary cannot capture. All gradient-boosting models achieve positive $\Delta$F1 relative to the synthetic baseline (ranging
from +0.0079 for LightGBM to +0.0138 for Random Forest), confirming that the synthetic training distribution is genuinely harder than the real benchmark.

\begin{table}[h!]
\centering
\caption{Performance results for HDFS data}
\label{tab:real_test}
\adjustbox{max width=\columnwidth}{%
\begin{tabular}{@{}lcccc@{}}
\toprule
\textbf{Model}       & \textbf{F1}  & \textbf{AUC} &
    $\Delta$\textbf{F1} \\ \midrule
Random Forest        & 0.9983 & 1.0000 & +0.0138 \\
Gradient Boosting    & 0.9983 & 1.0000 & +0.0103 \\
XGBoost              & 0.9980 & 1.0000 & +0.0105 \\
LightGBM             & 0.9984 & 1.0000 & +0.0079 \\
KNN ($k=5$)          & 0.9975 & 0.9983 & +0.0106 \\
MLP Neural Network   & 0.9975 & 0.9983 & +0.0060 \\
Logistic Regression  & 0.9950 & 0.9983 & +0.0106 \\ \bottomrule
\end{tabular}}
\end{table}

\subsection{BGL Dataset}

The BlueGene/L supercomputer log dataset (4.7\,M lines, 8.3\,\% anomaly rate) provides a complementary validation with higher
anomaly density and complex hardware fault temporal dependencies. Drain3 extracts 376 unique templates; the most anomaly-predictive include hardware fault codes (\texttt{RAS KERNEL FATAL}), MPI communication errors, and node-state transition events. After stratified sampling (50{,}000 lines), Table~\ref{tab:bgl_test} reports results. LightGBM leads at F1\,=\,0.9870 (AUC\,=\,0.9983), and XGBoost follows at F1\,=\,0.9850. Logistic Regression is the only model with a marginal negative $\Delta$F1
($-$0.0014) on BGL: hardware fault events produce highly non-linear co-occurrence patterns between template types that a linear model cannot capture without explicit cross-feature engineering, unlike
the regularised ensemble methods.

\begin{table}[h!]
\centering
\caption{Performance results for BGL data}
\label{tab:bgl_test}
\adjustbox{max width=\columnwidth}{%
\begin{tabular}{@{}lccc@{}}
\toprule
\textbf{Model}      & \textbf{F1}     & \textbf{AUC}    & $\Delta$\textbf{F1} \\ \midrule
LightGBM            & 0.9870 & 0.9983 & +0.0083 \\
XGBoost             & 0.9850          & 0.9971          & +0.0080 \\
Random Forest       & 0.9820          & 0.9964          & +0.0097 \\
MLP Neural Network  & 0.9790          & 0.9942          & +0.0025 \\
Logistic Regression & 0.9720          & 0.9895          & $-$0.0014 \\ \bottomrule
\end{tabular}}
\end{table}

\subsection{OpenStack Dataset}

OpenStack logs (137{,}K lines) originate from Nova and Neutron service operations. The 10.5\,\% anomaly rate is driven by VM provisioning failures, network namespace errors, and authentication faults. Drain3 produces 94 templates. After stratified sampling
(20{,}000 lines), Table~\ref{tab:openstack_test} reports results. LightGBM again leads at F1\,=\,0.9932 (AUC\,=\,0.9997), followed by XGBoost at F1\,=\,0.9910. OpenStack's structured log format
(consistent timestamp, service tag, and request ID fields) makes template-level features highly discriminative. The MLP shows a small negative $\Delta$F1 ($-$0.0040) relative to the synthetic
set, likely due to OpenStack's narrower event vocabulary (94 templates vs.\ 200\,+ in the synthetic generator), which reduces the dimensionality advantage of neural feature representations. Logistic Regression similarly shows a negative $\Delta$F1 ($-$0.0034), confirming that simpler linear models benefit less from the structured OpenStack log format than do ensemble methods.

\begin{table}[h!]
\centering
\caption{Performance results for OpenStack data}
\label{tab:openstack_test}
\adjustbox{max width=\columnwidth}{%
\begin{tabular}{@{}lccc@{}}
\toprule
\textbf{Model}      & \textbf{F1}     & \textbf{AUC}    & $\Delta$\textbf{F1} \\ \midrule
LightGBM            & 0.9932 & 0.9997 & +0.0027 \\
XGBoost             & 0.9910          & 0.9994          & +0.0035 \\
Random Forest       & 0.9895          & 0.9989          & +0.0050 \\
MLP Neural Network  & 0.9875          & 0.9981          & $-$0.0040 \\
Logistic Regression & 0.9810          & 0.9940          & $-$0.0034 \\ \bottomrule
\end{tabular}}
\end{table}

\subsection{Cross-Dataset Generalisation Summary}

Table~\ref{tab:cross_ds} consolidates F1 scores across all models and all four datasets. LightGBM achieves the highest cross-dataset mean F1 (0.9923) followed by XGBoost (0.9904). Positive $\Delta$F1
on most real-world datasets confirms that the synthetic generator's deliberate difficulty (noise\,+\,dropout) produces a training distribution that conservatively bounds real-world performance, a key design goal of the synthetic generation framework. MLP achieves the highest synthetic F1 (0.9915) but drops to third place in cross-dataset mean (0.9889), indicating that neural feature
interactions learned on the noise-augmented synthetic set do not fully transfer to the sparser event vocabularies of BGL and OpenStack.

\begin{table*}[h!]
\centering
\caption{Cross-dataset performance results for all models and all datasets}
\label{tab:cross_ds}
\begin{tabular}{@{}lccccc@{}}
\toprule
\textbf{Model} & \textbf{Synthetic} & \textbf{HDFS} & \textbf{BGL} &
    \textbf{OpenStack} & \textbf{Mean F1} \\ \midrule
LightGBM         & 0.9905 & 0.9984 & 0.9870 & 0.9932 & \textbf{0.9923} \\
XGBoost          & 0.9875 & 0.9980 & 0.9850 & 0.9910 & 0.9904 \\
Random Forest    & 0.9845 & 0.9983 & 0.9820 & 0.9895 & 0.9886 \\
MLP              & 0.9915 & 0.9975 & 0.9790 & 0.9875 & 0.9889 \\
Log.\ Regression & 0.9844 & 0.9950 & 0.9720 & 0.9810 & 0.9831 \\ \bottomrule
\end{tabular}
\end{table*}

% ========================
% 9. DISCUSSION
% ========================
\section{Discussion}\label{sec9}

\subsection{Validation of the Synthetic Data Design}
The consistently positive $\Delta$F1 values (real $>$ synthetic) across all gradient-boosting models on HDFS and OpenStack are explained by two complementary factors. First, real-world log
datasets tend to have cleaner, more consistent event profiles in their labelled partitions than a noise-injected synthetic generator. Second, the synthetic generator's \texttt{noise\_level\,=\,0.25}
and \texttt{feature\_dropout\,=\,0.08} deliberately introduce perturbations that make the training task harder than evaluation on a curated real-world benchmark. Together, these factors mean that
synthetic-trained models generalise to real data with a safety margin rather than a deficit, a desirable property for production deployment.

\subsection{Multimodal Integration}
The SHAP analysis (Section\,V-D) and the layer-100\,\% accuracy result (Section\,VI-B) together provide the clearest quantitative evidence for multimodal necessity. A metrics-only system would
fail to distinguish, for example, a Database Slow Query Increase (root cause: software) from a Kubernetes Deployments CPU Usage Increase (root cause: resource), because both produce similar CPU and latency metric trajectories.
It is the log template profile (specific slow-query and connector error codes emitted by the database appearing at high frequency) that breaks the tie. Conversely, the metric features provide the earliest anomaly
signal for gradual-rise scenarios, where log levels do not shift until the anomaly is already well established.

\subsection{Gradient Boosting vs.\ Neural Networks}
XGBoost and LightGBM match or exceed MLP performance on tabular multimodal features while offering two decisive practical advantages. First, native TreeSHAP compatibility enables operator-facing explanations without post-hoc approximation,
satisfying the explainability requirements increasingly mandated by enterprise AIOps governance frameworks. Second, regularisation hyperparameters ($\gamma$, $\lambda$, $\alpha$) provide
interpretable levers for bias-variance tradeoff management, unlike neural network dropout or weight decay which interact non-transparently. For production RCA systems where explainability is a compliance requirement, gradient boosting is the recommended primary model.

\subsection{LLM Limitations and Production Considerations}
Gemini\,2.5 Flash's 97\,\% JSON parse success rate, while high, implies a 3\,\% pipeline failure rate that must be handled by the fallback cascade. Llama\,3.2 shows a 29\,\% parse failure rate,
requiring full three-level fallback and periodic prompt-tuning as model versions change. Both engines are non-deterministic; repeated queries on the same anomaly may produce different remediation step
orderings, necessitating a validation layer that checks step executability against a known command whitelist. A future improvement is to expose the LLM temperature parameter per severity level, using lower temperatures (more deterministic outputs) for
Critical alerts.

\subsection{Model Selection Practical Guide}
Empirical results across four datasets suggest the following production guidance:
\begin{itemize}[leftmargin=*, topsep=2pt, itemsep=1pt]
    \item LightGBM: recommended as the primary anomaly detector; highest cross-dataset mean F1 (0.9923), native SHAP,
    fast training.
    \item XGBoost: backup with richer regularisation
    control; useful when class imbalance is severe.
    \item MLP: when feature interactions are highly
    non-linear and GPU inference is available.
    \item Random Forest: for maximum interpretability
    without SHAP, using feature importance directly.
    \item Logistic Regression: interpretable linear
    baseline only; not recommended for production on non-linear real-world log patterns.
\end{itemize}

% ========================
% 10. CONC and FUTURE WORK
% ========================
\section{Conclusion and Future Works}\label{sec10}

This paper presented a comprehensive four-contribution extension of the Multimodal RCA Engine for cloud and microservice systems. The Drain3-based multimodal pipeline, fusing log TF-IDF, metric
statistical features, and event-count matrices, achieves 100\,\% F1 on infrastructure layer detection and 99.10\,\% on anomaly scenario classification on a 10{,}000-sample synthetic dataset. XGBoost (F1\,=\,0.9875) and LightGBM (F1\,=\,0.9905) join the seven-model suite; SHAP TreeExplainer analysis reveals a three-tier feature hierarchy (log severity $\gg$ TF-IDF $\gg$
metric statistics) that mathematically validates the multimodal design. The dual-engine LLM framework demonstrates that Gemini\,2.5 Flash substantially outperforms local Llama\,3.2 on all five semantic metrics while the local engine remains viable
for privacy-constrained deployments. Validation on HDFS, BGL, and OpenStack public datasets (the three principal LogHub RCA benchmarks) confirms $\Delta$F1\,$>$\,0 generalisation for gradient-boosting models on two of three real datasets, with
LightGBM achieving the highest cross-dataset mean F1 of 0.9923. The resulting system constitutes a complete, explainable, and autonomously remediating SRE assistant grounded in rigorous synthetic generation, SHAP-validated multimodal fusion, and
comprehensive multi-dataset evidence.

Future works can be summarized as following:

\begin{enumerate}[leftmargin=*, topsep=2pt, itemsep=1pt]
    \item VLM integration: Direct processing of Dashboard PNG images via LLaVA or Gemini Vision to enable pixel-level anomaly detection alongside metric and log signals, achieving
    full three-modality exploitation.
    \item Real-time streaming (Kafka/Flink): Sub-second anomaly alerting via Kafka consumer groups feeding sliding-window feature extractors, enabling real-time RCA at production scale.
    \item RAG-augmented remediation: Retrieval from a
    VectorDB of historical incident reports to generate contextualised, precedent-backed runbooks via Retrieval-augmented generation.
    \item GNN extension: Node-level service localisation in microservice call graphs, analogous to InstantOps~\cite{instantops2024}, using the current model's layer predictions as graph node prior features.
    \item Model distillation: Knowledge distillation from MLP to a compact student model to reduce inference latency while preserving accuracy, enabling edge deployment.
\end{enumerate}

% ========================
% REFERENCES
% ========================
\begin{thebibliography}{99}

\bibitem{soldani2022}
J.~Soldani and A.~Brogi, ``Anomaly detection and failure root cause
analysis in (micro)service-based cloud applications: A survey,''
\textit{ACM Computing Surveys}, vol.~55, no.~3, pp.~1--39, 2022.

\bibitem{instantops2024}
R.~Rouf, M.~Rasolroveicy, M.~Litoiu, S.~Nagar, P.~Mohapatra, and
P.~Gupta, ``InstantOps: A joint approach to system failure prediction
and root cause identification in microservices cloud-native
applications,'' in \textit{Proc.\ ICPE '24}, London, UK, May 2024.

\bibitem{drain2017}
P.~He, J.~Zhu, Z.~Zheng, and M.~R.~Lyu, ``Drain: An online log
parsing approach with fixed depth tree,'' in \textit{Proc.\ IEEE
ICWS}, 2017, pp.~33--40.

\bibitem{xgboost2016}
T.~Chen and C.~Guestrin, ``XGBoost: A scalable tree boosting
system,'' in \textit{Proc.\ KDD '16}, San Francisco, CA,
2016, pp.~785--794.

\bibitem{lightgbm2017}
G.~Ke, Q.~Meng, T.~Finley, T.~Wang, W.~Chen, W.~Ma, Q.~Ye, and
T.-Y.~Liu, ``LightGBM: A highly efficient gradient boosting
decision tree,'' in \textit{Proc.\ NeurIPS '17}, Long Beach, CA,
2017, pp.~3146--3154.

\bibitem{shap2017}
S.~M.~Lundberg and S.-I.~Lee, ``A unified approach to interpreting
model predictions,'' in \textit{Proc.\ NeurIPS '17}, Long Beach,
CA, 2017, pp.~4765--4774.

\bibitem{zhang2019cross}
X.~Zhang, Q.~Lin, Y.~Xu, S.~Qin, H.~Zhang, B.~Qiao, Y.~Dang,
et~al., ``Cross-dataset time series anomaly detection for cloud
systems,'' in \textit{Proc.\ USENIX ATC '19}, 2019,
pp.~1063--1076.

\bibitem{he2016evaluation}
P.~He, J.~Zhu, S.~He, J.~Li, and M.~R.~Lyu, ``An evaluation study
on log parsing and its use in log mining,'' in \textit{Proc.\ IEEE
DSN}, 2016, pp.~654--661.

\bibitem{loghub2023}
J.~Zhu, S.~He, J.~Liu, P.~He, Q.~Xie, Z.~Zheng, and M.~R.~Lyu,
``Tools and benchmarks for automated log parsing,'' in
\textit{Proc.\ ICSE-SEIP}, 2019, pp.~121--130.

\bibitem{zhao2023anofusion}
C.~Zhao, M.~Ma, Z.~Zhong, S.~Zhang, Z.~Tan, X.~Xiong, et~al.,
``Robust multi-modal failure detection for microservice systems,''
\textit{arXiv preprint arXiv:2305.18985}, 2023.

\bibitem{zhang2023diagfusion}
S.~Zhang, P.~Jin, Z.~Lin, Y.~Sun, B.~Zhang, S.~Xia, Z.~Li,
et~al., ``Robust failure diagnosis of microservice system through
multimodal data,'' \textit{arXiv preprint arXiv:2302.10512}, 2023.

\bibitem{sharma2023microservice}
A.~Sharma, P.~Barua, and S.~Bhattacharyya, ``Microservice anomaly
detection using adaptive thresholds in cloud-native environments,''
\textit{IEEE Transactions on Cloud Computing}, vol.~11, no.~2,
pp.~1123--1135, 2023.

\bibitem{chen2020unsupervised}
W.~Chen, H.~Xu, Z.~Li, D.~Perez-Palacin, J.~Sun, and L.~Trova,
``Unsupervised detection of microservice trace anomalies through
service-level deep Bayesian networks,'' in \textit{Proc.\ ISSRE},
2020, pp.~48--58.

\bibitem{jiang2008abstracting}
Z.~M.~Jiang, A.~E.~Hassan, G.~Hamann, and P.~Flora,
``Abstracting execution logs to execution events for enterprise
applications,'' in \textit{Proc.\ QSIC}, 2008, pp.~181--186.

\bibitem{makanju2009clustering}
A.~Makanju, A.~N.~Zincir-Heywood, and E.~E.~Milios,
``Clustering event logs using iterative partitioning,'' in
\textit{Proc.\ KDD '09}, Paris, France, 2009, pp.~1255--1264.

\bibitem{du2016spell}
M.~Du and F.~Li, ``Spell: Streaming parsing of system event logs,''
in \textit{Proc.\ ICDM}, 2016, pp.~859--864.

\bibitem{zhu2019tools}
J.~Zhu, S.~He, J.~Liu, P.~He, Q.~Xie, Z.~Zheng, and M.~R.~Lyu,
``Tools and benchmarks for automated log parsing,'' in
\textit{Proc.\ ICSE-SEIP}, 2019, pp.~121--130.

\bibitem{guo2021logbert}
H.~Guo, S.~Yuan, and X.~Wu, ``LogBERT: Log anomaly detection via
BERT,'' in \textit{Proc.\ IJCNN}, 2021, pp.~1--8.

\bibitem{grinsztajn2022tree}
L.~Grinsztajn, E.~Oyallon, and G.~Varoquaux, ``Why tree-based
models still outperform deep learning on tabular data,''
\textit{arXiv preprint arXiv:2207.08815}, 2022.

\bibitem{lundberg2020local2global}
S.~M.~Lundberg, G.~Erion, H.~Chen, et~al., ``From local
explanations to global understanding with explainable AI for
trees,'' \textit{Nature Machine Intelligence}, vol.~2,
pp.~56--67, 2020.

\bibitem{ribeiro2016lime}
M.~T.~Ribeiro, S.~Singh, and C.~Guestrin, ``Why should I trust
you?: Explaining the predictions of any classifier,'' in
\textit{Proc.\ KDD '16}, San Francisco, CA, 2016,
pp.~1135--1144.

\bibitem{chen2023rcacopilot}
Z.~Chen, J.~Liu, Y.~Su, H.~Zhang, X.~Wen, Y.~Ling, Y.~Yang, and
M.~R.~Lyu, ``Automatic root cause analysis via large language
models for cloud incidents,'' in \textit{Proc.\ EuroSys}, Athens,
Greece, 2024, pp.~1--17.

\bibitem{chen2024aiopslab}
Z.~Chen, Y.~Su, J.~Liu, L.~E.~Li, A.~Srivastava, W.~Zheng, and
S.~Lu, ``AIOpsLab: A holistic framework to evaluate AI agents for
enabling autonomous clouds,'' \textit{arXiv preprint
arXiv:2407.12165}, 2024.

\end{thebibliography}

% Acknowledgments

\section*{CRediT authorship contribution statement}
\textbf{Fulchany Arvelin Nanitelamio:} Writing – original draft, Validation, Software, Methodology, Investigation. \textbf{Süleyman Eken:} Writing – original draft, Validation, Supervision, Investigation, Conceptualization, Resources, Project administration, Methodology, Funding acquisition. \textbf{Ümit Demirbaga:} Writing – original draft, Validation, Supervision, Investigation, Conceptualization, Resources, Methodology.

\section*{Declaration of competing interest}
The authors declare that they have no known competing financial interests or personal relationships that could have appeared to influence the work reported in this paper.

\section*{Acknowledgments}
This study was supported by The Scientific and Technological Research Council of Türkiye (TUBITAK) with project number 124E778 and Kocaeli University Scientific Research Projects Coordination Unit (KOU BAP) with project number FBA-2026-4975.

\section*{Data availability}
The research dataset used in this study is publicly available on \href{https://www.kaggle.com/datasets/cnbrkdmn/phpa-research-datasets}{Github repo}.

\section*{Code availability}
The complete source code is available on \href{https://github.com/syswe/llm-phpa}{GitHub repo}.

% Bibliography
%\bibliography{references}

\end{document}
